\documentclass[12pt,a4paper]{article}

\usepackage[utf8]{inputenc}
\usepackage[T1]{fontenc}
\usepackage[brazil]{babel}
\usepackage{lmodern}
\usepackage{geometry}
\usepackage{amsmath,amssymb}
\usepackage{booktabs}
\usepackage{siunitx}
\usepackage{array}
\usepackage{graphicx}
\usepackage{caption}
\usepackage{hyperref}
\usepackage{enumitem}
\usepackage{xcolor}
\usepackage{setspace}
\usepackage{float}

\geometry{margin=2.5cm}
\setstretch{1.12}
\emergencystretch=2em
\hypersetup{colorlinks=true, linkcolor=black, urlcolor=blue}
\sisetup{output-decimal-marker={,}}
\graphicspath{{figuras/}}

\title{Cálculo do circuito de polarização de Transistor Bipolar de Junção para acionamento de relé\\(DRIVER RELÉ)}
\author{Curso de Formação de Professores - Parque da Ciência Newton Freire Maia\\Autores: Aron da Rocha Battistella e Marcos Rocha}
\date{16/04/2026}

\newcommand{\Vcc}{V_{CC}}
\newcommand{\Vbb}{V_{BB}}
\newcommand{\Vbe}{V_{BE}}
\newcommand{\Vcesat}{V_{CE(\mathrm{sat})}}
\newcommand{\Ic}{I_C}
\newcommand{\Ie}{I_E}
\newcommand{\Ib}{I_B}
\newcommand{\Il}{I_L}
\newcommand{\Rb}{R_B}
\newcommand{\Rl}{R_L}
\newcommand{\Rpd}{R_{PD}}

\begin{document}

\maketitle

\section*{Resumo}
\sloppy
\noindent Este documento apresenta, passo a passo, o dimensionamento dos componentes constituintes do circuito de \emph{driver} para acionar um relé. O exemplo específico aqui selecionado trata de um relé de uso geral alimentado em \SI{5}{V} usando um transistor NPN BC548 como chave. Para tanto, a análise de circuitos necessária deve levar o transistor à condição de saturação, em que a corrente no coletor do transistor deve ser suficiente para alimentar a bobina do relé. São mostrados dois níveis de análise: uma dedução algébrica geral, útil para discutir a dependência da resistência de base com o ganho de corrente do transistor, e uma aplicação numérica com valores reais de componentes. Ao final, também se discute o papel do diodo de proteção 1N4007 ligado em paralelo com a bobina e a presença do resistor de \emph{pull-down}.

\section{Introdução}

Inúmeras aplicações em eletricidade e eletrônica exigem o acionamento de cargas cujas potências são incompatíveis com a disponibilidade de circuitos lógicos e de controle. Dessa forma, o acionamento de tais cargas, como por exemplo, eletrodomésticos, motores de corrente alternada (AC) e corrente contínua (DC), servo-motores, bobinas solenoides, entre outros, por meio de circuitos de automação, exige uma interface de acionamento, também chamada de \emph{driver}.

Um circuito \emph{driver} é capaz de acionar cargas de mais alta potência, usando uma fonte de alimentação própria, a partir do controle feito por um circuito de mais baixa potência, alimentado por uma segunda fonte de alimentação. Em montagens didáticas com Arduino, por exemplo, essa interface é necessária porque o pino digital não deve alimentar diretamente uma carga indutiva, como a bobina de um relé.

\section{Símbolos utilizados}

Antes da dedução, convém explicitar os símbolos usados ao longo do texto:

\begin{center}
\begin{tabular}{>{\centering\arraybackslash}p{2.5cm} p{10.0cm}}
\toprule
Símbolo & Significado \\
\midrule
$\Vcc$ & tensão da fonte que alimenta a carga no coletor, isto é, a bobina do relé \\
$\Vbb$ & tensão de comando aplicada à malha da base, por exemplo a tensão de saída de um pino digital \\
$\Vbe$ & queda de tensão entre base e emissor do transistor em condução \\
$\Vcesat$ & queda de tensão entre coletor e emissor quando o transistor está saturado \\
$\Ic$ & corrente de coletor do transistor \\
$\Ie$ & corrente de emissor do transistor \\
$\Ib$ & corrente de base do transistor \\
$\Il$ & corrente da bobina do relé, isto é, a corrente da carga \\
$\Rb$ & resistor de base, ligado entre o pino de controle e a base do transistor \\
$\Rl$ & resistência elétrica equivalente da bobina do relé \\
$\Rpd$ & resistor de \emph{pull-down}, ligado entre a base do transistor e o terra \\
$\beta$ & ganho de corrente contínua do transistor, definido pela razão aproximada entre $\Ic$ e $\Ib$ na região ativa \\
$V_{\text{relé}}$ & tensão efetiva aplicada sobre a bobina do relé \\
$V_{\Rb}$ & tensão elétrica no resistor de base \\
$V_D$ & queda de tensão direta no diodo de proteção quando ele conduz \\
$V_C$ & tensão no coletor do transistor \\
$V_L$ & tensão induzida na bobina do relé \\
$L$ & indutância da bobina do relé \\
$i$ & corrente instantânea na bobina \\
$t$ & tempo \\
\bottomrule
\end{tabular}
\end{center}

\section{Objetivo e hipótese de trabalho}

Considere o circuito clássico em que um pino digital aciona a base de um transistor BC548 por meio de um resistor de base, enquanto a bobina de um relé é ligada ao coletor do transistor. Neste contexto, o transistor opera como chave eletrônica: quando entra em condução, deve permitir a passagem da corrente exigida pela bobina; quando entra em corte, interrompe essa corrente.

Nesta nota, a corrente de interesse na carga é a corrente de coletor do transistor. No problema de chaveamento do relé, a condição relevante é
\begin{equation}
\Ic = \Il,
\label{eq:ic_il}
\end{equation}
onde $\Ic$ é a corrente de coletor e $\Il$ é a corrente exigida pela bobina do relé. A corrente de emissor $\Ie$ é ligeiramente maior, pois
\begin{equation}
\Ie = \Ic + \Ib,
\end{equation}
onde $\Ib$ é a corrente de base. Para o dimensionamento da carga no coletor, entretanto, a igualdade importante é a da corrente de coletor com a corrente da bobina.

\begin{figure}[htbp]
    \centering
    \includegraphics[width=0.43\textwidth]{../../img/malha_coletor.png}
    \caption{Malha do coletor. Imagem feita pelos autores.}
    \label{fig:malha_coletor}
\end{figure}

\section{Equações gerais da malha do coletor}

Se a bobina do relé pode ser modelada, em primeira aproximação, por uma resistência elétrica $\Rl$, então sua corrente nominal é dada por
\begin{equation}
\Il = \frac{V_{\text{relé}}}{\Rl}.
\label{eq:il}
\end{equation}

Na malha do coletor, uma forma mais detalhada é considerar a queda de tensão coletor-emissor do transistor. Nesse caso,
\begin{equation}
\Il = \frac{\Vcc - \Vcesat}{\Rl}.
\label{eq:il_vce}
\end{equation}

Quando o transistor está saturado, $\Vcesat$ costuma ser pequeno em comparação com a tensão de alimentação da bobina. Por isso, em uma primeira estimativa didática, é comum aproximar
\begin{equation}
V_{\text{relé}} \approx \Vcc.
\end{equation}

\section{Relação geral entre corrente de coletor e corrente de base}

No modelo de ganho de corrente contínua do transistor bipolar, vale a relação
\begin{equation}
\Ic = \beta\Ib,
\label{eq:ic_beta_ib}
\end{equation}
onde $\beta$ representa o ganho de corrente do transistor.

Isolando a corrente de base,
\begin{equation}
\Ib = \frac{\Ic}{\beta}.
\label{eq:ib}
\end{equation}

Como a corrente do coletor deve igualar a corrente da bobina, substituindo a Eq.~\eqref{eq:ic_il} em \eqref{eq:ib}, obtém-se
\begin{equation}
\Ib = \frac{\Il}{\beta}.
\label{eq:ib_il_beta}
\end{equation}

\section{Equação geral da resistência de base}

A malha de entrada, isto é, a malha da base, pode ser aproximada por uma fonte de tensão $\Vbb$, uma junção base-emissor com queda de tensão $\Vbe$ e o resistor de base $\Rb$.

\begin{figure}[htbp]
    \centering
    \includegraphics[width=0.38\textwidth]{../../img/malha_base.png}
    \caption{Malha da base. Imagem feita pelos autores.}
    \label{fig:malha_base}
\end{figure}

Aplicando a Lei de Kirchhoff para tensões na malha da base, temos
\begin{equation}
-\Vbb + V_{\Rb} + \Vbe = 0,
\label{eq:kirchhoff_base}
\end{equation}
onde $V_{\Rb}$ é a tensão no resistor de base.

Pela Lei de Ohm,
\begin{equation}
V_{\Rb} = \Ib\Rb.
\end{equation}

Portanto,
\begin{equation}
\Ib = \frac{\Vbb - \Vbe}{\Rb}.
\label{eq:ib_ohm}
\end{equation}

Isolando $\Rb$,
\begin{equation}
\Rb = \frac{\Vbb - \Vbe}{\Ib}.
\label{eq:rb_geral}
\end{equation}

Substituindo \eqref{eq:ib_il_beta} em \eqref{eq:rb_geral}, chega-se à expressão geral:
\begin{equation}
\boxed{\Rb = \frac{\beta\,(\Vbb - \Vbe)}{\Il}}.
\label{eq:rb_beta_il}
\end{equation}

Usando a expressão geral da corrente da bobina, Eq.~\eqref{eq:il}, também se pode escrever
\begin{equation}
\boxed{\Rb = \beta\,(\Vbb - \Vbe)\,\frac{\Rl}{V_{\text{relé}}}}.
\label{eq:rb_beta_rl}
\end{equation}

As Eqs.~\eqref{eq:rb_beta_il} e \eqref{eq:rb_beta_rl} são gerais: até este ponto, nenhum valor numérico específico de componente foi necessário.


\section{Saturação e saturação forte}

Quando o transistor é usado como chave para acionar um relé, o objetivo não é fazê-lo operar na região ativa, mas sim levá-lo à saturação. Na região ativa, a corrente de coletor é aproximadamente controlada pela corrente de base por meio da relação $\Ic = \beta\Ib$. Na saturação, por outro lado, o transistor já conduz intensamente, e a corrente de coletor passa a ser limitada principalmente pela carga externa, isto é, pela bobina do relé e pela fonte de alimentação.

Diz-se que há \textbf{saturação forte} quando a corrente de base é escolhida com margem suficiente para manter o transistor saturado mesmo quando as condições reais se afastam das condições ideais usadas no cálculo. Essa margem é importante porque o ganho $\beta$ não é um número fixo: ele varia de um transistor para outro, depende da corrente de coletor, depende da tensão coletor-emissor e também pode mudar com a temperatura.

Além disso, o circuito real está sujeito a variações ambientais e práticas: aquecimento do transistor durante o funcionamento, aquecimento da bobina do relé, pequenas variações na tensão de alimentação, tolerância dos resistores, tolerância da resistência da bobina e diferenças entre lotes de componentes. Em mudanças bruscas de temperatura, por exemplo, a queda $\Vbe$ e o ganho efetivo do transistor podem mudar. Se o projeto for feito sem margem, essas variações podem tirar o transistor da saturação, aumentando a queda $\Vcesat$, elevando a dissipação de potência no transistor e reduzindo a tensão efetiva aplicada à bobina do relé.

Neste documento, usa-se primeiro \textbf{100\% do ganho típico adotado} como cálculo inicial. Para o exemplo considerado, toma-se $\beta = 200$. Em seguida, acrescenta-se um fator de segurança equivalente a usar \textbf{50\% menos ganho} no cálculo, isto é,
\begin{equation}
\beta_{\text{projeto}} = 0{,}50 \cdot 200 = 100.
\end{equation}

Usar um ganho de projeto menor aumenta a corrente de base calculada e, consequentemente, leva à escolha de um resistor de base menor. Essa escolha favorece a saturação forte, pois torna o acionamento menos sensível a variações de temperatura, tolerâncias dos componentes, variações de alimentação e mudanças nas características reais do transistor.

A Figura~\ref{fig:curvas_bc548} complementa essa discussão ao mostrar, de forma didática, uma família de curvas de saída do transistor BC548 para diferentes valores de corrente de base. Nela, podem-se identificar visualmente as regiões de corte, região ativa e saturação, bem como o ponto de operação estimado quando se utiliza $\Rb = \SI{10}{k\ohm}$.

\begin{figure}[H]
    \centering
    \includegraphics[width=0.92\textwidth]{../../img/grafico_bc548_datasheet_curvas_ib_rb10k.png}
    \caption{Família de curvas de saída didáticas do transistor BC548 para diferentes valores de corrente de base, com indicação das regiões de corte, ativa e de saturação, da linha de carga do relé de \SI{5}{V} com bobina de \SI{70}{\ohm}, e do ponto de operação estimado para $\Rb = \SI{10}{k\ohm}$. Elaboração dos autores com base em valores do datasheet do BC548.}
    \label{fig:curvas_bc548}
\end{figure}

\section{Aplicação numérica com valores reais dos componentes}

A partir daqui, usam-se valores típicos da montagem didática. Considere:
\begin{equation}
\Vbb = \SI{5}{V},
\qquad
\Vcc = \SI{5}{V},
\qquad
\Vbe \approx \SI{0,7}{V},
\qquad
\Rl = \SI{70}{\ohm}.
\end{equation}

Como a bobina do relé está associada a uma alimentação de aproximadamente \SI{5}{V}, desprezando inicialmente a pequena queda $\Vcesat$, tem-se
\begin{equation}
V_{\text{relé}} \approx \SI{5}{V}.
\end{equation}

Logo, para o relé considerado,
\begin{equation}
\Il = \frac{V_{\text{relé}}}{\Rl}
= \frac{5}{70}
= \SI{0,0714}{A}
\approx \SI{71,4}{mA}.
\label{eq:il_srd_70}
\end{equation}

Para um pino digital de \SI{5}{V}, com $\Vbb = \SI{5}{V}$ e $\Vbe \approx \SI{0,7}{V}$, a Eq.~\eqref{eq:rb_beta_il} se torna
\begin{equation}
\Rb = \frac{\beta\,(5 - 0,7)}{\Il}
= \frac{4,3\,\beta}{\Il}.
\label{eq:rb_43}
\end{equation}


\subsection{Cálculo inicial usando 100\% do ganho típico: \texorpdfstring{$\beta = 200$}{beta = 200}}

Como cálculo inicial, adota-se o ganho típico de referência completo. Neste exemplo, considera-se
\begin{equation}
\beta = 200.
\end{equation}

Nesse caso, a corrente de base estimada é
\begin{equation}
\Ib = \frac{\Il}{\beta}
= \frac{71,4\,\mathrm{mA}}{200}
\approx \SI{0,357}{mA}.
\end{equation}

Logo,
\begin{equation}
\Rb = \frac{4,3}{0,000357}
\approx \SI{12045}{\ohm}.
\end{equation}

Portanto,
\begin{equation}
\boxed{\Rb \approx \SI{12,0}{k\ohm}}.
\end{equation}

\subsection{Cálculo com fator de segurança de 50\% no ganho: \texorpdfstring{$\beta = 100$}{beta = 100}}

Para favorecer a saturação forte, acrescenta-se uma margem de segurança usando um ganho de projeto 50\% menor que o ganho típico adotado. Assim,
\begin{equation}
\beta_{\text{projeto}} = 100.
\end{equation}

Com esse ganho reduzido, a corrente de base estimada é
\begin{equation}
\Ib = \frac{\Il}{\beta_{\text{projeto}}}
= \frac{71,4\,\mathrm{mA}}{100}
\approx \SI{0,714}{mA}.
\end{equation}

Logo,
\begin{equation}
\Rb = \frac{4,3}{0,000714}
\approx \SI{6022}{\ohm}.
\end{equation}

Portanto,
\begin{equation}
\boxed{\Rb \approx \SI{6,0}{k\ohm}}.
\end{equation}

Comparando os dois resultados, nota-se que o cálculo com $\beta = 100$ exige maior corrente de base e, portanto, menor resistência de base. Essa redução de $\Rb$ aumenta a margem de acionamento do transistor e ajuda a preservar a saturação caso ocorram variações bruscas de temperatura ou outras mudanças nas condições reais de operação.

\section{Diodo de proteção em paralelo com a bobina}

Como a bobina do relé é um elemento indutivo, ela se opõe a variações bruscas de corrente. Essa propriedade pode ser escrita como
\begin{equation}
V_L = L\frac{di}{dt}.
\end{equation}

Quando o transistor desliga, a corrente na bobina tende a cair rapidamente. Nesse instante, o termo $di/dt$ pode assumir grande valor em módulo, produzindo uma sobretensão capaz de danificar o transistor.

Para evitar esse problema, utiliza-se um diodo de proteção, por exemplo o 1N4007, em paralelo com a bobina e polarizado reversamente durante a operação normal. Assim, enquanto o relé está energizado, o diodo permanece desligado e não interfere no funcionamento do circuito.

No instante do desligamento, a polaridade induzida na bobina se inverte, o diodo entra em condução e fornece um caminho para a corrente de recirculação. Dessa forma, a tensão induzida deixa de crescer livremente e passa a ficar limitada aproximadamente à queda direta no diodo, da ordem de
\begin{equation}
V_D \approx \SI{0,7}{V}.
\end{equation}

Consequentemente, a tensão no coletor do transistor fica limitada a um valor aproximado de
\begin{equation}
V_C \approx \Vcc + V_D.
\end{equation}

Para $\Vcc = \SI{5,0}{V}$,
\begin{equation}
V_C \approx 5,0 + 0,7 = \SI{5,7}{V}.
\end{equation}

\section{Influência do diodo no cálculo da resistência de base}

Embora o diodo seja essencial para a proteção do transistor, ele \textbf{não altera o cálculo da resistência de base} durante a energização normal do relé. Isso ocorre porque, enquanto a bobina está alimentada, o diodo permanece reversamente polarizado e não participa da malha da base.

Assim, a equação de entrada continua sendo
\begin{equation}
\Vbb - \Vbe = \Rb\Ib,
\end{equation}
e, portanto,
\begin{equation}
\Rb = \frac{\Vbb - \Vbe}{\Ib}.
\end{equation}

Logo, a presença do diodo 1N4007 não modifica a dedução de $\Rb$. O diodo atua principalmente no instante de desligamento do relé, protegendo o transistor contra a sobretensão gerada pela bobina.

\section{Resistor de pull-down}

Além do resistor de base, é comum inserir um resistor de \emph{pull-down}, indicado por $\Rpd$, entre a base do transistor e o terra. Sua função é garantir que a base fique em nível lógico baixo quando o pino de controle estiver em alta impedância, desconectado ou durante instantes de inicialização do microcontrolador. Assim, evita-se que ruídos ou estados flutuantes acionem o relé indevidamente.

O resistor de \emph{pull-down} não deve competir de modo significativo com a corrente de base. Por isso, escolhe-se normalmente $\Rpd$ bem maior que $\Rb$, por exemplo na faixa de dezenas de quilohms. Nessa condição, a maior parte da corrente fornecida pelo pino de controle atravessa a junção base-emissor, e não o resistor de \emph{pull-down}.

\begin{figure}[H]
    \centering
    \includegraphics[width=0.80\textwidth]{../../img/diagrama_driver_rele.png}
    \caption{Diagrama esquemático do circuito \emph{driver} para acionamento do relé. A imagem representa a topologia geral do circuito, com transistor NPN, resistor de base, resistor de \emph{pull-down}, diodo de proteção e bobina do relé. Imagem feita pelos autores.}
    \label{fig:diagrama_driver_rele}
\end{figure}

\section{Resumo final}

A dedução geral deste documento é:
\begin{equation}
\boxed{\Rb = \frac{\beta\,(\Vbb - \Vbe)}{\Il}}
\qquad \text{com} \qquad
\Il = \frac{V_{\text{relé}}}{\Rl}.
\end{equation}

Considerando os valores reais adotados neste exemplo, $\Vbb = \SI{5}{V}$, $\Vbe \approx \SI{0,7}{V}$, $V_{\text{relé}} \approx \SI{5}{V}$ e $\Rl = \SI{70}{\ohm}$, obtém-se:

\begin{center}
\begin{tabular}{>{\raggedright\arraybackslash}p{5.0cm} >{\centering\arraybackslash}p{3.0cm} >{\centering\arraybackslash}p{3.6cm}}
\toprule
Hipótese de cálculo & Corrente de base & Resistência de base \\
\midrule
$\beta = 200$ (100\% do ganho típico) & \SI{0,357}{mA} & \SI{12,0}{k\ohm} \\
$\beta = 100$ (ganho 50\% menor) & \SI{0,714}{mA} & \SI{6,0}{k\ohm} \\
\bottomrule
\end{tabular}
\end{center}

\bigskip

\noindent\textbf{Conclusão prática:} o cálculo com $\beta = 200$ representa o uso de 100\% do ganho típico adotado como estimativa inicial. Já o cálculo com $\beta = 100$ introduz um fator de segurança de 50\% menos ganho, aumentando a corrente de base prevista e favorecendo a saturação forte. Assim, a escolha prática de $\Rb$ deve considerar não apenas a conta ideal, mas também variações de temperatura, tolerâncias dos componentes, variações de alimentação e o limite de corrente seguro do pino de controle.

Para a proteção do transistor contra sobretensões no desligamento, considera-se também um diodo 1N4007 em paralelo com a bobina do relé, polarizado reversamente durante o funcionamento normal. O resistor de \emph{pull-down}, por sua vez, ajuda a evitar acionamentos indesejados quando a entrada de controle estiver flutuante.

\section*{Referências iniciais}

\begin{enumerate}[leftmargin=1.0cm]
    \item BOYLESTAD, Robert L.; NASHELSKY, Louis. \textit{Dispositivos eletrônicos e teoria de circuitos}. 11. ed. São Paulo: Pearson, 2013.
    \item ON Semiconductor. \textit{BC546/547/548/549/550 - 30 V, 100 mA NPN general-purpose transistors}. Datasheet.
    \item SONGLE Relay. \textit{SRD series relay datasheet}. Dados típicos de bobina para a família SRD-05VDC-SL-C.
\end{enumerate}

\end{document}
